\documentclass[a4paper,12pt]{article}
\usepackage[utf8]{inputenc}
\usepackage[spanish]{babel}
\usepackage{amsmath}
\usepackage{amsfonts}
\usepackage{amssymb}
\usepackage{siunitx}
\usepackage{geometry}
\usepackage{enumitem}

\geometry{margin=2.5cm}

\title{Examen de Física I - Forma C}
\author{Nombre: \underline{\hspace{8cm}} \quad Matrícula: \underline{\hspace{3cm}}}
\date{Fecha: \underline{\hspace{4cm}}}

\begin{document}

\maketitle

\section{Parte I: Teoría (4 puntos - 1 pt c/u)}
Responda brevemente a las siguientes preguntas:

\begin{enumerate}[itemsep=2em]
    \item \textbf{Definición}: ¿Qué estudia la ciencia de la Física?
    \item \textbf{Redondeo (Prod/Div)}: ¿Cuál es la regla para determinar las cifras significativas en una multiplicación o división?
    \item \textbf{Magnitudes Fundamentales}: Enumere cinco (5) magnitudes fundamentales del Sistema Internacional (SI) junto con sus unidades básicas.
    \item \textbf{Notación Científica}: ¿Cómo se expresaría el número 0.00000045 en notación científica? Explique brevemente la regla del exponente negativo.
\end{enumerate}

\vspace{2cm}

\section{Parte II: Práctica (6 puntos - 1.5 pts c/u)}
Resuelva los siguientes problemas mostrando todo su procedimiento. Redondee según las reglas de cifras significativas aprendidas.

\begin{enumerate}[itemsep=4em]
    \item \textbf{Conversión de Longitud}: Convierta \SI{2.5}{\km} a pulgadas (in). (Use $1\text{ in} = 2.54\text{ cm}$).
    \item \textbf{Conversión de Caudal}: Convierta \SI{1.5}{\gallon\per\second} a $\text{cm}^3/\text{hr}$. (Use \SI{1}{\gallon} = \SI{3785.41}{\centi\metre^3}).
    \item \textbf{Notación Científica}: Realice la siguiente operación y exprese el resultado en notación científica: $(6.0 \times 10^{-4}) \times (2.5 \times 10^{-5})$.
    \item \textbf{Cifras Significativas}: Realice la siguiente multiplicación y redondee correctamente: $12.45 \times 2.0$.
\end{enumerate}

\end{document}
